\par\textbf{Desarrollo para comparar tu intento.} Para comparar poblaciones con distinta exposición, calcular eventos/horas en cada una. El recuento absoluto por sí solo omite la duración de observación.

\begin{enumerate}

\item Normalizar los cuatro eventos por su exposición.

\[\fitmath{\mathrm{rateA}=\frac{4}{20000}\approx 0.0002\,\mathrm{1/h}}\]

\item Normalizar los seis eventos por la otra exposición.

\[\fitmath{\mathrm{rateB}=\frac{6}{60000}\approx 0.0001\,\mathrm{1/h}}\]

\item Comparar tasas observadas; el cociente no prueba causalidad ni significación estadística.

\[\fitmath{\mathrm{ratio}=\frac{\frac{6}{60000}}{\frac{4}{20000}}\approx 0.5\,}\]

\end{enumerate}
